\subsection*{Preface}

What is new here is the degree to which our typology both creates and helps us
to discover the self and its defenses in a personal way. It is thus a
contribution to the understanding of the person in the psyche. That person is
urged into being by archetypes, including the transpersonal Self that seems to
drive human individuation; but the 'little-s' self we experience as persons has
its own forms of consciousness, nested in the care of archetypes but capable of
asserting and integrating themselves independently. I say this lest anyone
reading what I have written about the archetypes in this book take the
eight-function, eight-archetype model as an unaltering image of inevitable
rigidity. Rather, it is the psyche's many dynamic parts that give it the
flexibility, over time, to articulate the Self in a personal way that our
adaptation to our typology makes us capable of displaying. The reader is
therefore advised to approach the book with this developmental perspective in
mind. It may not be amiss, as I celebrate in these pages a full fifty years of
reflecting on Jungian typology, to admit what by now must be obvious: that my
own perspective on the subject is continuing to develop.

John Beebe

San Francisco, October 2015
